\cotefile{63}
\scanpage{36}
\[
\boxed{y^2+a_1xy+a_3y=x^3+a_2x^2+a_4x+a_6}
\]
\[
\beta_2=a_1^2+4a_2\qquad \beta_4=a_1a_3+2a_4\qquad \beta_6=a_3^2+4a_6
\]
\[
-\beta_8=a_1a_3a_4+a_4^2-a_1^2a_6-a_2a_3^2-4a_2a_6
\]
\[
\boxed{\unsure{-4}\beta_8=\beta_4^2-\beta_2\beta_6}
\]
\[
\inserted{\beta_8=\beta_2a_6+a_2a_3^2-a_1a_3a_4\deleted{\unsure{+}}-a_4^2
=a_2\beta_6+a_1^2a_6-a_4^2-a_1a_3a_4}
\]
\margin{$\unsure{g_2},\quad\unsure{g_3}$}
\[
\gamma_4=\beta_2^2-24\beta_4\qquad -\gamma_6=\beta_2^3-36\beta_2\beta_4+216\beta_6
\]
\[
\begin{aligned}
\Delta&=-\beta_2^2\beta_8-8\beta_4^3-27\beta_6^2+9\beta_2\beta_4\beta_6\\
&=\beta_4^3-27\beta_6^2+\beta_8(36\beta_4-\beta_2^2)
\end{aligned}
\qquad\boxed{1728\,\Delta=\gamma_4^3-\gamma_6^2}
\]
If $C':\quad y'^{\,2}+a_1'x'y'+\cdots$ is another curve in the same form, then a birational transformation $C'\xrightarrow{\varphi}C$ keeping $O\ (x=\infty,y=\infty)$ fixed will be of the form
\margin{$\deleted{\unsure{x\circ\varphi}}=$}
\[
\left\{\begin{aligned}
x\circ\varphi&=\rho^2x'+r\\
y\circ\varphi&=\rho^3y'+s\rho^2x'+t
\end{aligned}\right.\qquad \rho\ne0
\]
\margin{\tikz[baseline=(v.base)]{\node (v) {$j=\dfrac{\gamma_4^3}{\Delta}$};\draw (v.center) ellipse (0.95cm and 0.65cm);}}
The coefficients are then related by
\[
\begin{aligned}
\rho a_1'&=a_1+2s\\
\rho^2a_2'&=a_2-sa_1+3r-s^2\\
\rho^3a_3'&=a_3+ra_1+2t\\
\rho^4a_4'&=a_4-sa_3+2ra_2-(t\deleted{\unsure{-}}+rs)a_1+3r^2-2st\\
\rho^6a_6'&=a_6+ra_4+r^2a_2+r^3-ta_3-t_{\deleted{\unsure{a_1}}}^2-rta_1\ !
\end{aligned}
\]
\margin{$-\dfrac{\partial F}{\partial y}(r,t)$\\
$\dfrac{\partial F}{\partial x}(r,t)+\unsure{s|a}\dfrac{\partial F}{\partial y}(r,t)\unsure{=}$\\
$F(r,t)=$}
\margin{$\downarrow\quad\unsure{-45}\quad\unsure{-60}$}
\margin{$\deleted{\unsure{rs}}$}
\margin{\tikz[baseline=(v.base)]{\node (v) {$\begin{aligned}12g_2&=\gamma_4\\216g_3&=+\gamma_6\end{aligned}$};\draw (v.center) ellipse (1.35cm and 0.67cm);}}
The transformation of the $\beta$'s is
\[
\begin{aligned}
\rho^2\beta_2'&=\beta_2+12r\\
\rho^4\beta_4'&=\beta_4+r\beta_2+6r^2\\
\rho^6\beta_6'&=\beta_6+2r\beta_4+r^2\beta_2+4r^3\\
\rho^8\beta_8'&=\beta_8+3r\beta_6+\deleted{\unsure{2}}3r^2\beta_4+r^3\beta_2+3r^4
\end{aligned}
\]
\margin{$\mathrm{disc}\unsure{=|-}\gamma_4$\\
$\mathrm{disc}=\overset{?}{16}\Delta$}
\margin{\tikz[baseline=(v.base)]{\node (v) {$\begin{aligned}j\Delta&=\gamma_4^3\\(j-1728)\Delta&=\gamma_6^2\end{aligned}$};\draw (v.center) ellipse (1.82cm and 0.68cm);}}
From these we check that
\[
\rho^4\gamma_4'=\gamma_4\qquad\rho^6\gamma_6'=\gamma_6
\]
\margin{\tikz[baseline=(v.base)]{\node (v) {$\dfrac{j}{j-1728}=\dfrac{\gamma_4^3}{\gamma_6^2}$};\draw[thick] (v.center) ellipse (1.48cm and 0.65cm);}}
And now using $1728\,\Delta=\gamma_4^3-\gamma_6^2$ we see
\[
\rho^{12}\Delta'=\Delta
\]
\[
\deleted{\unsure{\rho\omega}}\ \omega'=\rho\omega=\frac{dx}{2y+a_1x+a_3}
\]
\margin{$\unsure{\omega\circ\varphi|\omega\cdot\varphi}=\rho^{-1}\omega'$ \ (see over)}
\[
\boxed{\boxed{\begin{aligned}
y^2+xy&=x^3-\frac{36\deleted{\unsure{j}}}{j-1728}\,x
-\frac{\deleted{\unsure{j}}1}{j-1728}
&\text{\textquotedblleft generic }A_j\text{\textquotedblright}\\
\gamma_4&=\gamma_6=\frac{j}{j-1728}\ ;\qquad
\Delta=\frac{j^2}{(j-1728)^3}
&j\ne0,1728,\unsure{\infty}
\end{aligned}}}
\]
% archive sheet number: 35

\scanpage{37}
\mspage{2}
\[
F=y^2+a_1xy+a_3y-x^3-a_2x^2-a_4x-a_6
\]
\[
F_y=2y+a_1x+a_3\qquad -F_x=3x^2+2a_2x+a_4-a_1y
\]
The point at $\infty$ is non-singular \unsure{so|as} the singularities are the common zeros of $F,F_x,F_y$. Given $a_1,a_2,a_3,a_4,a_6$ we get a curve of genus 1 if and only if there are no singularities, and in that case, a differential of first kind is
\[
\omega=\frac{dx}{2y+a_1x+a_3}=\frac{dy}{3x^2+2a_2x+a_4-a_1y}
\]
% margin note written sideways
\margin{\raisebox{0pt}[0pt][0pt]{\rotatebox{-90}{\scriptsize$\omega=\dfrac{dx}{2y+a_1x+a_3}
=\dfrac{\rho^2dx'}{2\rho^3y'+(2s\rho^2x'+a_1\rho^2x')+(2t+a_1r+a_3)}
=\dfrac1{\rho}\dfrac{dx'}{2y'+a_1'x'+a_3'}=\dfrac1{\rho}\omega'$}}}
It transforms according to the rule
\[
\omega'=\rho\omega
\]
\uline{Special} \uline{rational} \uline{canonical} \uline{forms}:
\[
\text{\uuline{$p\ne2,3$}}\qquad
\boxed{y^2=x^3+a_4x+a_6}\qquad\omega=\frac{dx}{2y}
\]
\[
\gamma_4=-48a_4,\quad\gamma_6=216\cdot4a_6,\quad
\Delta=16(-4a_4^3-27a_6^2)
\]
Singularity $=$ double root of cubic $f(x)=x^3+a_4x+a_6$. $\Delta$ is 16 times the discriminant of $f(x)$. Thus $g=1\Longleftrightarrow\Delta\ne0$. The allowable transformations are simply $x=\rho^2x',\quad y=\rho^3y'$.
\[
j=\frac{\gamma_4^3}{\Delta}
=1728\,\frac{4a_4^3}{4a_4^3+27a_6^2}
=1728-1728\,\frac{27a_6^2}{4a_4^3+27a_6^2}
\]
is an absolute invariant. Given $C$ and $C'$ with same $j$, they are birationally equivalent over the algebraic closure $\bar{k}$ of $k$:

\uuline{case 1} \quad $j\ne0,1728$\qquad i.e. $a_4\ne0,\ a_6\ne0$.
\[
C\sim C'\Longleftrightarrow j=j'\ \text{and}\ 
\frac{a_4}{a_6}\frac{a_6'}{a_4'}\in k^{*^{2}}\ ;\quad
k'=k\left(\sqrt{\frac{a_4a_6'}{a_6a_4'}}\right)
\]
Automorphisms \quad $\rho=\pm1$. \qquad Canonical curve $C_j$ ?

\uuline{Case 2} \quad $j=1728$\qquad i.e. $a_6=0$
\[
C\sim C'\Longleftrightarrow j=j'\ \text{and}\ 
\frac{a_4}{a_4'}\in k^{*^{4}}\ ;\quad
k'=k\left(\sqrt[4]{\frac{a_4}{a_4'}}\right)
\]
Automorphisms \quad $\rho=\pm1$ and $\pm i$ if $i\in k$

Canonical curve $C_{1728}\ \deleted{\unsure{=}}\quad y^2=x^3-x$

\uuline{Case 3} \quad $j=0$\qquad i.e. $a_4=0$
\[
C\sim C'\Longleftrightarrow j'=j\ \text{and}\ 
\frac{a_6}{a_6'}\in k^{*^{6}}\ ;\quad
k'=k\left(\sqrt[6]{\frac{a_6}{a_6'}}\right)
\]
Automorphisms \quad $\rho^6=1$, \quad $\rho=\pm1$, \quad and more if $k$ contains $\unsure{\sqrt{-3}|1\!:\!3}$

Canonical curve $C_0:\qquad y^2=x^3-1$
% archive sheet number: 36

\scanpage{38}
\mspage{3}
\uuline{$p=3$}\qquad Preliminary. \quad $y^2=x^3+a_2x^2+a_4x+a_6=f(x)$
\[
\begin{aligned}
\beta_2&=a_2 &\beta_4&=-a_4 &\beta_8&=a_4^2-a_2a_6\\
\gamma_4&=a_2^2 &\gamma_6&=a_2^3 &\Delta&=a_2^2a_4^2-a_2^3a_6\deleted{\unsure{+}}-a_4^3
\end{aligned}
\]
Again, $\Delta=\mathrm{disc}\,f(x)$ so $g=1\Longleftrightarrow\Delta\ne0$.
\[
j=\frac{a_2^6}{\Delta}
\]
\uuline{Case 1}: \quad $j\ne0$ i.e. $a_2\ne0$.
\[
\boxed{y^2=x^3+a_2x^2+a_6}\qquad \omega=\frac{dx}{2y}
\]
\[
\Delta=-a_2^3a_6\qquad j=-\frac{a_2^3}{a_6}
\]
Allowable transformations: \quad $x=\rho^2x',\ y=\rho^3y'$.
\[
C=C'\Longleftrightarrow j=j'\quad\text{and}\quad
\frac{a_2}{a_2'}\in k^{*^{2}}\qquad k'=k\left(\sqrt{\frac{a_2}{a_2'}}\right)
\]
Automorphisms \qquad $\rho=\pm1$

Canonical curve ?

\uuline{Case 2} \quad $j=0$ i.e. $a_2=0$
\[
\boxed{y^2=x^3+a_4x+a_6}
\]
$\Delta=-a_4^3$, so $a_4\ne0$

allowable transformations: \qquad $x=\rho^2x'+r\qquad y=\rho^3y'$
\[
\rho^4a_4'=a_4\qquad \rho^6a_6'=a_6+ra_4+r^3
\]
$C=C'\Longleftrightarrow j=j'=0$\quad and \quad $\dfrac{a_4'}{a_4}\in k^{*^{4}}$ \quad and $\left(\dfrac{a_4}{a_4'}\right)^{\frac32}a_6'-a_6$ of form $r^3+ra_4$

\deleted{\unsure{Automs} \quad $\rho=\pm1,\pm i\quad r^3+a_4r+\unsure{(a_6\pm a_6')} $}

$k'/k$ separable of degree $\le12$
\[
\begin{aligned}
\text{Automs}\quad \rho&=\pm1 & r^3+a_4r+(a_6-a_6')&=0\\
\rho&=\pm i & r^3+a_4r+(a_6+a_6')&=0
\end{aligned}
\]
\margin{group of order \uuline{12}}
Canonical curve $C_0$ \unsure{:} \quad $y^2=x^3-x$. Here $\rho^4=1\quad r^3-r=0$.

$G_{12}$.
% archive sheet number: 37

\scanpage{39}
\mspage{4}
\uuline{$p=2$}\qquad $\rho a_1'=a_1$

\uuline{Case 1}: \quad $a_1\ne0$: \quad $a_1=1$ fixes $\rho=1$; $a_3=0$ fixes $r=0$, $a_4=0$ fixes $t=0$
\[
\boxed{y^2+xy=x^3+a_2x^2+a_6}\qquad \omega=\frac{dx}{x}
\]
\[
\beta_2=1\qquad \beta_4=\beta_6=0\qquad\beta_8=a_6\qquad\gamma_4=1\qquad\Delta=a_6
\]
\[
F_y=x,\quad F_x=3x^2+y\qquad \text{singularity}\Longleftrightarrow\Delta=0
\qquad j=\frac1{a_6}
\]
\[
\begin{aligned}
a_2'&=a_2+s^2-s\\
a_6'&=a_6
\end{aligned}
\]
\[
C\sim C'\Longleftrightarrow j=j'\ \text{and}\quad
a_2'-a_2\in\unsure{\wp|p}(k)\qquad k'=k\bigl(\wp^{-1}(a_2'-a_2)\bigr)
\]
Automs \quad $s=0,1$

Canonical curve \quad $y^2+xy=x^3+a_6$

\uuline{Case 2}: \quad $a_1=0$ \qquad can make $a_2=0$\qquad fixes $r=s^2$
\[
\boxed{y^2+a_3y=x^3+a_4x+a_6}\qquad \omega=\frac{dx}{a_3}
\]
\[
\begin{aligned}
\rho^3a_3'&=a_3\\
\rho^4a_4'&=a_4+sa_3+s^4\\
\rho^6a_6'&=a_6+s^2a_4+ta_3+s^6+t^2
\end{aligned}
\quad\left|\quad
\begin{aligned}
&\beta_2=\beta_4=0,\ \beta_6=a_3^2,\ \beta_8=a_4^2\\
&\Delta=a_3^4\qquad\gamma_4=0\qquad j=0\\
&F_y=a_3\qquad F_x=x^2+a_4\\
&\text{non sing}\Longleftrightarrow\Delta\ne0
\end{aligned}\right.
\]
$C\sim C'\Longleftrightarrow j=j'=0$ and $\dfrac{a_3'}{a_3}\in k^3$ and the other two equations are solvable. \quad $k'/k$ is separable of degree $\le24$

Canonical Curve: \quad $y^2-y=x^3$

Automs of it: \qquad $\rho^3=1\qquad s^4+s=0\qquad t^2+t+s^3+s^2=0$

$G$ of order 24:
% archive sheet number: 38

\scanpage{40}
\mspage{5}
\uline{Summarizing}: If $j\ne0,1728$ (any characteristic) then in a given field $k$ there are as many distinct curves of invariant $j$ as there are separable quadratic extensions (with 1 additional) and each of these has only the reflection \unsure{autom} (and the trivial one). Given two such curves they determine a unique quad ext over which they are equivalent

In \uline{all} cases two curves with the same $j$ are birat equivalent over a \uline{separable} extension of degree dividing 24, and have an automorphism group which is a subgroup of the group of order 24.
% archive sheet number: 39

\scanpage{66}
\[
x_2,y_3,g_4,g_6\qquad \boxed{y_3^2=x_2^3+g_4x_2+g_6}
\]
\[
\left\{\begin{aligned}
x&=x_{\omega}=\omega^{-2}x_2\\
y&=y_{\omega}=\omega^{-3}y_3\\
\deleted{g_4}\gamma_4&=\gamma_{4\omega}=\omega^{-4}g_4\\
\gamma_6&=\gamma_{6\omega}=\omega^{-6}g_6
\end{aligned}\right.
\qquad\text{\uline{i.e.}}\quad y^2=x^3+\gamma_4x+\gamma_6
\]
\[
\boxed{dx=2y\omega}\quad\text{i.e.}\quad
\bigl\lfloor\underline{d(\omega^{-2}x_2)=2\omega^{-2}y_3}\bigr\rfloor
\]
\[
2y\,dy=(3x^2+\gamma_4)dx=2y(3x^2+\gamma_4)\omega
\]
\[
\boxed{dy=\deleted{\unsure{\tfrac12}}(3x^2+\gamma_4)\omega}\qquad
\bigl\lfloor\underline{d(\omega^{-3}y_3)=\deleted{\unsure{\tfrac12}}(3\deleted{\unsure{\omega^{-4}}}x_2^2+
\deleted{\unsure{\omega^{-4}}}g_4)\omega^{-3}}\bigr\rfloor
\]
\[
\left\{\begin{aligned}
\omega\deleted{\unsure{y}}&=\frac12\frac{dx}{y}\\
\eta&=\omega x=\deleted{\unsure{dx}}\frac12\frac{x\,dx}{y}
\quad\overset{\unsure{\omega}}{\bigl[}=\omega^{-1}x_2\bigr]
\end{aligned}\right.
\]
\[
\left\{\begin{aligned}
\int_{\gamma_1}\omega&=\omega_1,\quad \omega_1=\int_0^{\omega_1}dz\\
\int_{\gamma_2}\omega&=\omega_2,\quad \omega_2=\int_0^{\omega_2}dz
\end{aligned}\right.
\qquad X^{\mathrm{an}}\simeq C/E(\omega_1,\omega_2)
\]
\[
\left\{\begin{aligned}
x&=\frac1{z^2}+\sum_{a\in E(\omega_1,\omega_2)}'\left[\frac1{(z-a)^2}-\frac1{a^2}\right]\\
y&=-\frac1{z^3}\unsure{-}\deleted{\unsure{2}}\sum_{a\in E(\omega_1,\omega_2)}\frac1{(z-a)^3}
\end{aligned}\right.
\]
\[
\left\{\begin{aligned}
\omega&=dz\\
\eta&=x\,dz
\end{aligned}\right.
\qquad\eta_1=\int_u^{u+\omega_1}x\,dz,\quad
\eta_2=\int_u^{u+\omega_2}x\,dz
\]
\[
\begin{aligned}
\gamma_4&=\\
\gamma_6&=
\end{aligned}
\]
% archive sheet number: 65

\scanpage{67}
Passage du splitting standard des formes différentielles de 2\textsuperscript{\unsure{ème|e}} espèce en $\omega,\eta$ au splitting \unsure{transcendant} (de Hodge) $\omega,\bar\omega$, donc on aura
\[
\left\{\begin{aligned}
\eta+\alpha\omega&=\beta\bar{\omega}\\
\alpha&=-\frac{\bar{\omega}_1\eta_2-\bar{\omega}_2\eta_1}
{\bar{\omega}_1\omega_2-\bar{\omega}_2\omega_1}
\end{aligned}\right.
\]
\margin{N.B. dénominateur $=2i\,|\omega_1|^2\,\mathrm{Im}\,\tau\ne0$}
\uline{NB}\quad On a une 2 forme \unsure{invariante|invariant} canonique associée à la comparaison des deux splittings
\[
\pi=\alpha\omega^2\quad\in\unsure{\Gamma}\operatorname{Hom}(\underline\omega^{-1},\underline\omega)=\Gamma\underline\omega^2
\]
\uline{Questions}\qquad a) Quand $\alpha=0$ ?, i.e. $\bar\omega$ proportionnel : $\eta$ ??
b) Quand α est-il
b) Quand $\alpha$ est-il algébrique sur $\unsure{\mathbb Q|Q}$ ?
% archive sheet number: 66

\scanpage{68}
\mspage{- 19 -}\hfill n\textsuperscript{o}\handwritten{357}

\uline{pour tout} $\handwritten{\lambda}\handwritten{\in}\handwritten{\mathscr{C}}'(M)$, $U(x)$ \uline{est scalairement essentiellement} $\handwritten{\lambda}$-\uline{intégrable au sens de la topologie de la convergence compacte dans} $\handwritten{\mathscr{L}}(E,E)$ ; \uline{on a} $U(\handwritten{\lambda})=\handwritten{\int} U d\handwritten{\lambda}\handwritten{\in}\handwritten{\mathscr{L}}(E,E)$, \uline{et}, \uline{quels que soient} $u\handwritten{\in} E$ \uline{et} $u'$ \uline{dans le dual} $E'$ \uline{de} $E$ :
\begingroup\makeatletter\tagsleft@true\makeatother
\[
U(\handwritten{\lambda})u=\handwritten{\int} U(x)u.d\handwritten{\lambda}(x)\quad,\quad
\handwritten{\langle} U(\handwritten{\lambda})u,u'\handwritten{\rangle}=\handwritten{\int}\handwritten{\langle} U(x)u,u'\handwritten{\rangle}\,d\handwritten{\lambda}(x)\quad;\tag*{(9)}
\]
\uline{de plus}, \uline{quels que soient} $\handwritten{\lambda}\handwritten{\in}\handwritten{\mathscr{C}}'(M)$, $\mu\handwritten{\in}\handwritten{\mathscr{C}}'(M)$, \uline{on a} :
\[
U(\handwritten{\lambda}\ast\mu)=U(\handwritten{\lambda})U(\mu)\ .\tag*{(10)}
\]
\endgroup
\uline{Si enfin l'ensemble des} $U(x)$, \uline{pour} $x\handwritten{\in} M$, \uline{est équicontinu}, \uline{les mêmes résultats subsistent lorsqu'on substitue} $\handwritten{\mathscr{M}}^{1}(M)$ \uline{à} $\handwritten{\mathscr{C}}'(M)$.

Le fait que $U(\handwritten{\lambda})$ est défini et est dans $\handwritten{\mathscr{L}}(E,E)$ est un cas particulier de la prop.16 du Chap.VI, \handwritten{\S}\,1,n\textsuperscript{o}7 ; la seconde des relations (9) résulte alors de la définition même de $\handwritten{\int} U d\handwritten{\lambda}$, puisque $\handwritten{\langle} Uu,u'\handwritten{\rangle}$ est une forme linéaire continue sur $\handwritten{\mathscr{L}}(E,E)$ (prop.11 de E.V.T., Chap.IV, \handwritten{\S}\,2, n\textsuperscript{o}9, joint au fait que la topologie de la convergence compacte est plus fine que celle de la convergence simple) [N.B. - Le rédacteur est tout fier d'avoir pu donner une raison aussi savante pour un fait aussi évident, mais n'a pas trouvé dans Bourbaki d'autre référence pour celui-ci]. La première suit de là en vertu des définitions. [N.B. - Les formules (9) auraient en réalité dû figurer au Chap.VI, \handwritten{\S}\,1, n\textsuperscript{o}7 !]. En substituant $\handwritten{\lambda}\ast\mu$ à $\handwritten{\lambda}$ dans la seconde formule (9), et appliquant la définition de $\handwritten{\lambda}\ast\mu$, on obtient :
\[
\handwritten{\langle} U(\handwritten{\lambda}\ast\mu)u,u'\handwritten{\rangle}
=\handwritten{\int}\handwritten{\langle} U(xy)u,u'\handwritten{\rangle}\,d\handwritten{\lambda}(x)d\mu(y)
=\handwritten{\int}\handwritten{\langle} U(y)u,{}^tU(x)u'\handwritten{\rangle}\,d\handwritten{\lambda}(x)d\mu(y)\ .
\]
Par hypothèse, $x\handwritten{\longrightarrow} U(x)u$ est continue, donc $\handwritten{\langle} U(xy)u,u'\handwritten{\rangle}$ est continue sur $M\times M$ ; comme l'intégrabilité essentielle d'une fonction continue entraîne son intégrabilité, on peut appliquer le théorème de Lebesgue-Fubini au deuxième membre, ou, ce qui revient au même au troisième membre de la relation ci-dessus ; cela donne, d'après la première relation (9) :
% archive sheet number: 67

\scanpage{69}
\[
\begin{aligned}
\eta_1&=\left[-\frac1{\unsure{z|3}}\right]_u^{u+\omega_1}
+\sum_{a\ne0}\bigl(\left[-\frac1{z-a}-\frac1{a^2}z\right]_u^{u+\omega}\\
&=\left(\frac1u-\frac1{u+\omega_1}\right)
+\sum_{a\ne0}\left[\frac1{u-a}-\frac1{u+\omega_1-a}-\frac{\omega_1}{a^2}\right]\\
&=\omega_1\left[\frac1{u(u+\omega_1)}+
\sum_{a\ne0}\left(\underbrace{\frac1{(u-a)(u+\omega_1-a)}-\frac1{a^2}}_{
\displaystyle\frac1{a^3}\frac{(2u+\omega_1)-u(u+\omega_1)a^{-1}}{(1-ua^{-1})(1-(u+\omega_1)a^{-1})}}\right)\right]
\end{aligned}
\]
\[
\deleted{=\omega}
\]
\[
=\omega_1\left[\underbrace{\frac1{u(u+\omega_1)}+\frac1{(u-\omega_1)u}}_{
\displaystyle\frac2{\deleted{\unsure{u}}(u+\omega_1)(u-\omega_1)}}
-\frac1{\omega_1^2}+\sum_{\substack{a\ne0\\ a\ne\omega_1}}
\quad\tikz[baseline=-2pt]{\draw[->] (0,.48) .. controls (.25,.53) and (.16,.19) .. (.36,.19) -- (.86,.22) -- (1.03,.02);}\qquad\right]
\]
Faisons $u\leadsto0$, \unsure{on trouve}
\[
\eta_1=\omega_1\left[-\frac3{\omega_1^2}
+\sum_{\substack{a\ne0\\a\ne\omega_1}}
\underbrace{\frac1{a^3}\frac{\omega_1}{1-\omega_1a^{-1}}}_{
\displaystyle\frac1{a^3}\frac1{\omega_1^{-1}-a^{-1}}}\right]
\]
\[
\eta_1=-\frac3{\omega_1}+
\omega_1\sum_{\substack{a\ne0\\a\ne\omega_1}}
\left(\underbrace{\frac1{a^3}\frac1{\omega_1^{-1}-a^{-1}}}_{
\displaystyle\frac1{a^2}\frac{\omega_1}{a-\omega_1}}\right)
=-\frac3{\omega_1}+\omega_1^2\sum_{\substack{a\ne0\\a\ne\omega_1}}
\frac1{a^2}\cdot\frac1{a-\omega_1}
\]
\[
a=m\omega_1+n\omega_1
\]
\[
\eta_2=\hspace{5em}-\frac3{\omega_2}+\omega_2^2\sum_{\substack{a\ne0\\a\ne\omega_1}}
\frac1{a^2}\frac1{a-\omega_2}
\]
\[
\deleted{\unsure{a\ne|a+}}
\]
% archive sheet number: 68

\scanpage{70}
\mspage{- 7 -}\hfill n\textsuperscript{o} \handwritten{380}

Dans le premier cas, on observe que $\lVert e-x_n^{-1}x\rVert=\lVert x_n^{-1}(x_n-x)\rVert\handwritten{\le}\lVert x_n^{-1}\rVert\ \lVert x_n-x\rVert$ donc tend vers zéro lorsque $n\handwritten{\to}\handwritten{\infty}$. Par suite, pour $n$ assez grand, $x_n^{-1}x$ est inversible, ce qui implique que $x$ est inversible.

Dans le second cas, compte tenu de la prop.2, e), on a $\handwritten{\unsure{\mathscr{D}|\mathfrak{D}|\partial}}(e-x_n^{-1}x)\handwritten{\le}\handwritten{\mathscr{D}}(x_n^{-1})\,\handwritten{\mathscr{D}}(x_n-x)\handwritten{\le}\handwritten{\mathscr{D}}(x_n^{-1})\ \lVert x_n-x\rVert$, donc $\handwritten{\mathscr{D}}(e-x_n^{-1}x)$ tend vers 0 lorsque $n\handwritten{\to}\handwritten{\infty}$, ce qui implique que pour $n$ assez grand $x_n^{-1}x$ est inversible et par conséquent, $x$ est inversible.
\begin{center}\texttt{-:-:-:-}\end{center}
% archive sheet number: 69

\scanpage{71}
\[
\begin{aligned}
\eta_1&=-\frac3{\omega_1^{\deleted{\unsure{1|2}}}}+\omega_1^2
\sum_{\substack{(m,n)\ne(0,0)\\(m,n)\ne(\deleted{\unsure{0}}1,0)}}
\frac1{(m\omega_1+n\omega_2)^2}\frac1{(m-1)\omega_1+n\omega_2}\\
\eta_2&=-\frac3{\omega_2}+\omega_2^2
\sum_{\substack{(m,n)\ne(0,0)\\(m,n)\ne(0,1)}}
\frac1{(m\omega_1+n\omega_2)^2}\frac1{(m\omega_1+(n-1)\omega_2)}
\end{aligned}
\]
\begingroup\scriptsize
\[
\begin{aligned}
\bar\omega_1\eta_2-\bar\omega_2\eta_1
&=-3\left(\frac{\bar\omega_1}{\omega_2}-\frac{\bar\omega_2}{\omega_1}\right)
\deleted{+\sum_{\substack{(m,n)\ne(0,0)\\(1,0)\\(0,1)}}\bar\omega_1\omega_2^2}\\
&\quad+\sum_{\scriptstyle (m,n)\ne\left|\substack{(0,0)\\(1,0)\\(0,1)}\right.}
\frac{|\omega_1|^2\bigl((m-1)\omega_2^2+n\frac{\omega_2^3}{\omega_1}\bigr)
-|\omega_2|^2\bigl((n-1)\omega_1^2+m\frac{\omega_1^3}{\omega_2}\bigr)}
{(m\omega_1+n\omega_2)^2\bigl((m-1)\omega_1+n\omega_2\bigr)\bigl(m\omega_1+(n-1)\omega_2\bigr)}\\
&\quad+\underbrace{\bar\omega_1\omega_2^2\frac1{\omega_1^2(\omega_1-\omega_2)}
-\bar\omega_2\omega_1^2\frac1{\omega_2^2(\omega_2-\omega_1)}}_{
\displaystyle\frac{\bar\omega_1\omega_2^4+\bar\omega_2\omega_1^4}
{(\omega_1\omega_2)^2(\omega_1-\omega_2)^{\deleted{\unsure{2}}}}}
\end{aligned}
\]
\endgroup
\[
=\frac{\bar\omega_1\omega_2^4+\bar\omega_2\omega_1^4}
{\deleted{(}\omega_1^{\unsure{1|2}}\omega_2^2(\omega_1-\omega_2)\deleted{\unsure{)^2|)}}}
-3\frac{|\omega_1|^2-|\omega_2|^2}{\omega_1\omega_2}
+\sum_{\scriptstyle(m,n)\ne\left(\substack{0,0\\1,0\\0,1}\right.}
\quad\tikz[baseline=-2pt]{\draw[->] (2.9,1.05) -- (2.9,.8) .. controls (2.9,.68) and (2.82,.05) .. (2.15,.02) -- (.02,0);}\,
\]
\begingroup\scriptsize
\[
=\frac{\deleted{\unsure{\bar\omega_1|\omega_1}}\lvert\omega_1\rvert^2}{\omega_1^2}
\frac{\tau^4+\bar\tau}{(\tau^2(1-\tau)\deleted{\unsure{)^2|)}}}
-\frac{|\omega_1|^2}{\omega_1^2}\frac{1-|\tau|^2}{\tau}
+\sum_{\scriptstyle(m,n)\ne\left|\substack{00\\01\\10}\right.}\frac{|\omega_1|^2}{\omega_1^2}
\frac{(m-1)\tau^2+n\tau^3-(n-1)|\tau|^2+m|\tau|^2\tau^{\unsure{-1|-}}}
{(m+n\tau)^2((m-1)+n\tau)(m+(n-1)\tau)}
\]
\endgroup
\begingroup\scriptsize
\[
=\frac{|\omega_1|^2}{\omega_1^2}\left(
\frac{\tau^5+|\deleted{\unsure{\omega}}\tau|^2}{\tau^3(1-\tau)}
-\frac{1-|\tau|^2}{\tau}
+\sum_{\scriptstyle(m,n)\ne\left|\substack{00\\01\\10}\right.}
\underbrace{\frac{n\tau^4+(m-1)\tau^3+(n-1)|\tau|^2\tau+m|\tau|^2}
{\tau(m+n\tau)^2(m-1+n\tau)(m+(n-1)\tau)}}\right)
\]
\endgroup
\begingroup\scriptsize
\[
\deleted{\bar\omega_1}\quad
\alpha=-\frac{\bar\omega_1\eta_2-\bar\omega_2\eta_1}{2i|\omega_1|^2\mathrm{Im}\,\tau}
=-\frac{1}{2i\omega_1^2\mathrm{Im}(\tau)}
\left(\frac{\tau^5+|\tau|^2}{\tau^3(1-\tau)}
-\frac{1-|\tau|^2}{\tau}
+\sum_{\scriptstyle(m,n)\ne\left|\substack{0,0\\0,1\\1,0\\\deleted{\unsure{0,1|0}}}\right.}
\quad\tikz[baseline=-1pt]{\draw[->] (0,.7) .. controls (.5,.7) and (.8,.45) .. (.73,.03);}\right)
\]
\endgroup
% archive sheet number: 70

\scanpage{72}
\mspage{- \handwritten{6} -}\hfill n\textsuperscript{o} \handwritten{380}

3. \uline{Inverses}

Dans ce N\textsuperscript{o} on suppose que l'algèbre $A$ est normée \uline{complète} et possède un élément unité $e$.

PROPOSITION 4.- \uline{Pour tout élément} $x\handwritten{\in}A$ \uline{tel que} $\handwritten{\unsure{\mathscr{D}|\mathfrak{D}|\partial}}(e-x)<1$, \uline{la série} $\displaystyle\handwritten{\sum}_{u=1}^{\handwritten{\infty}}(e-x)^n$ \uline{est absolument convergente et} $\displaystyle e+\handwritten{\sum}_{u=1}^{\handwritten{\infty}}(e-x)^n$ \uline{est l'inverse de} $x$.

La série est absolument convergente d'après le Critère de Cauchy. Soit $\displaystyle x'=e+\handwritten{\sum}_{u=1}^{\handwritten{\infty}}(e-x)^n$. On a
\[
x'x=xx'=x+\handwritten{\sum}_{u=1}^{\handwritten{\infty}}(e-x)^n x=x+\handwritten{\sum}_{u=1}^{\handwritten{\infty}}(e-x)^n-\handwritten{\sum}_{u=1}^{\handwritten{\infty}}(e-x)^{n+1}=e.
\]
COROLLAIRE.- \uline{Le groupe des éléments inversibles de} $A$ \uline{contient la boule ouverte de centre} $e$ \uline{de rayon} $1$.

En effet $\handwritten{\mathscr{D}}(x-e)\handwritten{\le}\lVert x-e\rVert$ pour tout $x$.

Ce résultat figurait en Top. Gén., Chap.9, \handwritten{\S} 3, prop.13. On a vu par la même occasion que le groupe des éléments inversibles de $A$ est un ouvert de $A$ et que la topologie induite sur ce groupe par la topologie de $A$ est compatible avec la structure de groupe. Si $A$ est de dimension finie sur $C$, ce groupe ne possède qu'un nombre fini de composantes connexes. Dans les autres cas, il arrive qu'il possède un nombre infini de composantes connexes.

PROPOSITION 5.- \uline{Soit} $A$ \unsure{'}\uline{une algèbre normée complète avec unité} $e$ \uline{et soit} $x_n$ \uline{une suite convergente d'éléments inversibles de} $A$ \uline{ayant pour limite} $x$ :

a) \uline{si la suite} $\lVert x_n^{-1}\rVert$ \uline{est bornée}, \uline{alors} $x$ \uline{est inversible},

a') \uline{si la suite} $\handwritten{\mathscr{D}}(x_n^{-1})$ \uline{est bornée et si} $xx_n=x_nx$ \uline{pour tout} $n$, \uline{alors} $x$ \uline{est inversible}.

\uline{Dans les deux cas}, \uline{l'inverse de} $x$ \uline{est la limite des} $x_n^{-1}$.
% archive sheet number: 71

